\chapter{The REPL Is Your Laboratory}\label{ch:repl}

\index{REPL}

A chemistry lab doesn't start with a production facility.
It starts with a bench, some glassware, and the freedom to mix things and see what happens.
The Lattice REPL is that bench.

The REPL---\textbf{R}ead, \textbf{E}valuate, \textbf{P}rint, \textbf{L}oop---is where you'll spend much of your time when learning Lattice.
It's where you test a hypothesis about how a function works, where you prototype a data structure before committing it to a file, and where you debug a tricky expression by poking at it from different angles.
In this chapter, we'll learn how to use the REPL effectively and discover the features that make it more than a toy calculator.

%% ───────────────────────────────────────────────────────────
\section{Launching the REPL}\label{sec:launching-repl}
\index{REPL!launching}

Start the REPL by running \texttt{clat} with no arguments:

\begin{lstlisting}[language=bash, numbers=none]
./clat
\end{lstlisting}

You'll see the welcome banner:

\begin{lstlisting}[numbers=none]
Lattice v0.3.28 -- crystallization-based programming language
Copyright (c) 2026 Alex Jokela. BSD 3-Clause License.
Type expressions to evaluate. Ctrl-D to exit.
\end{lstlisting}

The \lstinline|lattice>| prompt awaits your input.
Type an expression and press Enter:

\begin{lstlisting}[language=Lattice, numbers=none]
lattice> 6 * 7
=> 42
\end{lstlisting}

To exit, press \texttt{Ctrl-D} (which sends an EOF signal) or close the terminal.

\subsection{Persistent State Across Lines}\label{subsec:persistent-state}
\index{REPL!persistent state}

The REPL is not a series of disconnected calculations.
Every variable, function, struct, and enum you define persists for the rest of the session.
Think of it as a single, growing program that you write one line at a time:

\begin{lstlisting}[language=Lattice, numbers=none]
lattice> let city = "Portland"
lattice> let state = "Oregon"
lattice> print("${city}, ${state}")
Portland, Oregon
lattice> let population = 652_503
lattice> population
=> 652503
\end{lstlisting}

The variable \lstinline|city| defined on line one is still available on line three.
The variable \lstinline|population| defined on line four is available for the rest of the session.

This is possible because the REPL maintains a single, long-lived virtual machine instance.
Under the hood (in \filename{src/main.c}), each line you type is compiled to bytecode and executed on the same \lstinline|StackVM|---globals, functions, and type definitions are preserved between compilations.

\begin{defnbox}{REPL Persistence}
The REPL uses a \emph{persistent VM}\index{persistent VM}: a single \lstinline|StackVM| instance that survives across lines.
Variables stored as globals, function definitions, struct declarations, and enum types all remain accessible until you exit.
This is what allows you to build programs interactively, one piece at a time.
\end{defnbox}

\subsection{Multi-line Input}\label{subsec:multi-line}
\index{REPL!multi-line input}

What happens when you type an incomplete expression---like opening a brace without closing it?

\begin{lstlisting}[language=Lattice, numbers=none]
lattice> fn add(a: Int, b: Int) -> Int {
    ...>     a + b
    ...> }
lattice> add(3, 4)
=> 7
\end{lstlisting}

The REPL detects that the input is incomplete by checking whether brackets, braces, and parentheses are balanced.
When they are not, the prompt changes to \lstinline|    ...>| and waits for more input.
The accumulated text is only compiled and executed once the delimiters are balanced.

This works for any kind of nesting: function bodies, \lstinline|if|/\lstinline|else| blocks, \lstinline|match| expressions, struct definitions, and more.

\begin{lstlisting}[language=Lattice, numbers=none]
lattice> let message = if true {
    ...>     "yes"
    ...> } else {
    ...>     "no"
    ...> }
lattice> message
=> "yes"
\end{lstlisting}

\begin{tipbox}{Escaping a Multi-line Mistake}
If you start typing a multi-line expression and realize you've made a mistake, press \texttt{Ctrl-C} to cancel the current input and return to a fresh \lstinline|lattice>| prompt.
Alternatively, close all open delimiters to force evaluation (which will likely produce an error, but clears the buffer).
\end{tipbox}

%% ───────────────────────────────────────────────────────────
\section{Tab Completion and Discoverability}\label{sec:tab-completion}
\index{tab completion}\index{REPL!tab completion}

The Lattice REPL supports tab completion powered by libedit (or readline, depending on your system).
Press \texttt{Tab} and the REPL will suggest completions based on what you've typed so far.

\subsection{Keyword and Built-in Completion}\label{subsec:keyword-completion}
\index{tab completion!keywords}

Start typing a keyword or built-in function name and press \texttt{Tab}:

\begin{lstlisting}[numbers=none]
lattice> fr<Tab>
freeze
\end{lstlisting}

If there are multiple matches, pressing \texttt{Tab} twice shows all options:

\begin{lstlisting}[numbers=none]
lattice> to<Tab><Tab>
to_string  to_int  to_float  toml_parse  toml_stringify
\end{lstlisting}

The completion engine knows about all of Lattice's keywords (\lstinline|fn|, \lstinline|let|, \lstinline|flux|, \lstinline|fix|, \lstinline|match|, \lstinline|struct|, \lstinline|enum|, \lstinline|trait|, \lstinline|impl|, and more) and over 120 built-in functions.
This makes the REPL a discovery tool: if you vaguely remember that there's a function related to files, type \lstinline|file| and press \texttt{Tab}:

\begin{lstlisting}[numbers=none]
lattice> file<Tab><Tab>
file_exists  file_size
\end{lstlisting}

\subsection{Method Completion After a Dot}\label{subsec:method-completion}
\index{tab completion!methods}

The completion engine is context-aware.
When the cursor follows a dot (\lstinline|.|), it switches to \emph{method mode}\index{tab completion!method mode} and suggests method names instead of global keywords:

\begin{lstlisting}[numbers=none]
lattice> "hello".to<Tab><Tab>
to_upper  to_lower
\end{lstlisting}

\begin{lstlisting}[numbers=none]
lattice> [1, 2, 3].fi<Tab><Tab>
filter  find  first  flat  flat_map  fill
\end{lstlisting}

The completion source for methods includes array, string, map, set, channel, buffer, and enum methods---all drawn from the same list defined in \filename{src/completion.c}.

\subsection{Constructor Completion After ::}\label{subsec:constructor-completion}
\index{tab completion!constructors}

Type a type name followed by \lstinline|::| and press \texttt{Tab} to see available constructors\index{constructors}:

\begin{lstlisting}[numbers=none]
lattice> Map::<Tab><Tab>
new
lattice> Set::<Tab><Tab>
new  from
lattice> Channel::<Tab><Tab>
new
\end{lstlisting}

This triple context---global, method, and constructor---means the REPL's tab completion adapts to where you are in an expression.

\begin{notebox}{How Completion Works Internally}
The tab completion system is implemented in \filename{src/completion.c}.
It registers a custom \lstinline|rl_attempted_completion_function| with the readline library.
When you press \texttt{Tab}, this function inspects the character immediately before the cursor:
\begin{itemize}
  \item If it's a dot (\lstinline|.|), completion mode switches to \texttt{COMPLETE\_METHOD}.
  \item If it's \lstinline|::|, completion mode switches to \texttt{COMPLETE\_SCOPE}.
  \item Otherwise, it uses \texttt{COMPLETE\_GLOBAL}, which walks through keywords, built-in functions, and constructors.
\end{itemize}
Filename completion is explicitly disabled---the REPL suggests Lattice symbols, not files.
\end{notebox}

%% ───────────────────────────────────────────────────────────
\section{Defining Functions, Structs, and Enums Interactively}\label{sec:defining-interactively}
\index{REPL!defining functions}\index{REPL!defining structs}\index{REPL!defining enums}

The REPL is not limited to one-liners.
You can define complex types and build up a small program incrementally.

\subsection{Functions}\label{subsec:repl-functions}

Define a function and call it immediately:

\begin{lstlisting}[language=Lattice, numbers=none]
lattice> fn celsius_to_fahrenheit(c: Float) -> Float {
    ...>     c * 9.0 / 5.0 + 32.0
    ...> }
lattice> celsius_to_fahrenheit(100.0)
=> 212.0
lattice> celsius_to_fahrenheit(0.0)
=> 32.0
\end{lstlisting}

The function persists.
You can call it again later, pass it to higher-order functions, or use it inside other definitions:

\begin{lstlisting}[language=Lattice, numbers=none]
lattice> let temps_c = [0.0, 20.0, 37.0, 100.0]
lattice> temps_c.map(|c| celsius_to_fahrenheit(c))
=> [32.0, 68.0, 98.6, 212.0]
\end{lstlisting}

\subsection{Structs}\label{subsec:repl-structs}
\index{structs!in REPL}

Define a struct, create instances, and call their methods:

\begin{lstlisting}[language=Lattice, numbers=none]
lattice> struct Point {
    ...>     x: Float,
    ...>     y: Float
    ...> }
lattice> let origin = Point { x: 0.0, y: 0.0 }
lattice> let target = Point { x: 3.0, y: 4.0 }
lattice> target.x
=> 3.0
\end{lstlisting}

Add a trait and implement it:

\begin{lstlisting}[language=Lattice, numbers=none]
lattice> trait Measurable {
    ...>     fn distance(self: Any) -> Float
    ...> }
lattice> impl Measurable for Point {
    ...>     fn distance(self: Any) -> Float {
    ...>         sqrt(self.x * self.x + self.y * self.y)
    ...>     }
    ...> }
lattice> target.distance()
=> 5.0
\end{lstlisting}

\subsection{Enums}\label{subsec:repl-enums}
\index{enums!in REPL}

Enums work the same way:

\begin{lstlisting}[language=Lattice, numbers=none]
lattice> enum Direction {
    ...>     North,
    ...>     South,
    ...>     East,
    ...>     West
    ...> }
lattice> let heading = Direction::North
lattice> heading.variant_name()
=> "North"
lattice> heading.is_variant("North")
=> true
\end{lstlisting}

\begin{lstlisting}[language=Lattice, numbers=none]
lattice> enum Shape {
    ...>     Circle(Float),
    ...>     Rectangle(Float, Float)
    ...> }
lattice> let s = Shape::Rectangle(4.0, 5.0)
lattice> s.payload()
=> [4.0, 5.0]
\end{lstlisting}

The ability to define, test, and iterate on types interactively is one of the REPL's great strengths.
You can experiment with a struct's field layout or an enum's variants before committing them to a source file.

%% ───────────────────────────────────────────────────────────
\section{The \texttt{=>} Prefix: How the REPL Talks Back}\label{sec:repl-output}
\index{=> prefix (REPL)}\index{REPL!output conventions}

You've already seen the \lstinline|=>| prefix in front of REPL results.
Let's understand exactly when it appears and what it means.

\subsection{The Display Rule}\label{subsec:display-rule}

The REPL follows a clear rule: \textbf{if the last expression in your input produces a value that is not \lstinline|unit| and not \lstinline|nil|, it is displayed with the \lstinline|=>| prefix.}

\begin{lstlisting}[language=Lattice, numbers=none]
lattice> 42
=> 42
lattice> "hello"
=> "hello"
lattice> [1, 2, 3]
=> [1, 2, 3]
lattice> true
=> true
\end{lstlisting}

Statements that produce \lstinline|unit|\index{unit} are silently suppressed:

\begin{lstlisting}[language=Lattice, numbers=none]
lattice> let x = 5
lattice> flux counter = 0
lattice> counter += 1
lattice>
\end{lstlisting}

No \lstinline|=>| output.
Variable declarations and assignments return \lstinline|unit|, so there is nothing to display.

The \lstinline|print()| function writes directly to standard output---it does \emph{not} produce a REPL result:

\begin{lstlisting}[language=Lattice, numbers=none]
lattice> print("hello")
hello
lattice>
\end{lstlisting}

Notice: \lstinline|hello| appears without the \lstinline|=>| prefix.
That's \lstinline|print()| writing to stdout.
The \lstinline|print()| function itself returns \lstinline|unit|, so the REPL suppresses its return value.

\subsection{The \texttt{repr} Format}\label{subsec:repr-format}
\index{repr}

The REPL uses the \lstinline|repr|\index{repr} format for display, not \lstinline|to_string|.
The difference matters for strings:

\begin{lstlisting}[language=Lattice, numbers=none]
lattice> "hello world"
=> "hello world"
\end{lstlisting}

The quotes are part of the display---they tell you this value is a string.
If you \lstinline|print()| the same value, the quotes are absent:

\begin{lstlisting}[language=Lattice, numbers=none]
lattice> print("hello world")
hello world
\end{lstlisting}

This distinction between \lstinline|repr| (for programmers, with type information) and \lstinline|print| (for users, raw output) is consistent throughout Lattice.
Structs that define a custom \lstinline|repr| closure field will use it for REPL display.

\subsection{Distinguishing \texttt{print} from \texttt{=>}}\label{subsec:print-vs-arrow}

Here is a common source of confusion for newcomers.
What does this produce?

\begin{lstlisting}[language=Lattice, numbers=none]
lattice> print(2 + 2)
4
lattice>
\end{lstlisting}

The \lstinline|4| is from \lstinline|print()|.
The REPL has no \lstinline|=>| line because \lstinline|print()| returns \lstinline|unit|.
Compare with:

\begin{lstlisting}[language=Lattice, numbers=none]
lattice> 2 + 2
=> 4
\end{lstlisting}

Here the \lstinline|4| is the REPL displaying the expression result.
Understanding this distinction helps you debug REPL sessions: if you see \lstinline|=>|, it's the value of the expression.
If you don't, it's side-effect output (like \lstinline|print()|).

\begin{defnbox}{Unit and Nil}
\lstinline|unit|\index{unit} is the return type of statements and functions with no meaningful return value---like variable assignments or \lstinline|print()|.
\lstinline|nil|\index{nil} is an explicit null value that means ``no value here.''
Both are suppressed in REPL output, but they are different types:
\lstinline|typeof(nil)| returns \lstinline|"Nil"| and \lstinline|typeof(unit)| returns \lstinline|"Unit"|.
We will cover both in \cref{ch:values-types}.
\end{defnbox}

%% ───────────────────────────────────────────────────────────
\section{Tips and Tricks for Exploratory Programming}\label{sec:repl-tips}
\index{REPL!tips and tricks}

The REPL is more than an input box.
Here are techniques that will make your interactive sessions more productive.

\subsection{Use \texttt{typeof()} to Inspect Values}\label{subsec:typeof-repl}
\index{typeof()}

When you're not sure what type a value has, ask:

\begin{lstlisting}[language=Lattice, numbers=none]
lattice> typeof(42)
=> "Int"
lattice> typeof(3.14)
=> "Float"
lattice> typeof("hello")
=> "String"
lattice> typeof([1, 2, 3])
=> "Array"
lattice> typeof(Map::new())
=> "Map"
lattice> typeof(nil)
=> "Nil"
lattice> typeof(true)
=> "Bool"
\end{lstlisting}

This is especially useful when a function returns an unexpected result and you want to know what you're working with.

\subsection{Use \texttt{phase\_of()} to Inspect Phase}\label{subsec:phaseof-repl}
\index{phase\_of()}

Check whether a value is fluid or crystal:

\begin{lstlisting}[language=Lattice, numbers=none]
lattice> flux items = [1, 2, 3]
lattice> phase_of(items)
=> "fluid"
lattice> freeze(items)
lattice> phase_of(items)
=> "crystal"
\end{lstlisting}

\subsection{Build Data Structures Incrementally}\label{subsec:incremental-building}
\index{REPL!incremental building}

The REPL is perfect for building up complex data one step at a time:

\begin{lstlisting}[language=Lattice, numbers=none]
lattice> flux inventory = Map::new()
lattice> inventory["apples"] = 12
lattice> inventory["bananas"] = 6
lattice> inventory["oranges"] = 8
lattice> inventory
=> {"apples": 12, "bananas": 6, "oranges": 8}
lattice> inventory.keys()
=> ["apples", "bananas", "oranges"]
lattice> inventory.values().sum()
=> 26
\end{lstlisting}

You can see the state of your data at each step, catch mistakes early, and iterate quickly.

\subsection{Test Functions with Edge Cases}\label{subsec:edge-cases}

Once you've defined a function, throw edge cases at it:

\begin{lstlisting}[language=Lattice, numbers=none]
lattice> fn safe_divide(a: Float, b: Float) -> Float {
    ...>     if b == 0.0 { return 0.0 }
    ...>     a / b
    ...> }
lattice> safe_divide(10.0, 3.0)
=> 3.3333333333333335
lattice> safe_divide(10.0, 0.0)
=> 0.0
lattice> safe_divide(-7.5, 2.5)
=> -3.0
\end{lstlisting}

The REPL gives you a tight feedback loop for testing behavior before writing formal tests.

\subsection{Use History to Save Typing}\label{subsec:history}
\index{REPL!history}

The REPL supports command history via the readline library.
Press the \textbf{Up} and \textbf{Down} arrow keys to navigate through previous inputs.
This is handled by libedit/readline and works the same as in bash or Python's interactive interpreter.

Non-empty lines are automatically added to the history buffer, so you can recall and modify earlier definitions without retyping them.

\subsection{Explore the Standard Library}\label{subsec:explore-stdlib}
\index{REPL!exploring built-ins}

The REPL is the fastest way to explore Lattice's built-in functions.
Wondering what \lstinline|version()|\index{version()} returns?

\begin{lstlisting}[language=Lattice, numbers=none]
lattice> version()
=> "0.3.28"
\end{lstlisting}

Curious about the current platform?

\begin{lstlisting}[language=Lattice, numbers=none]
lattice> platform()
=> "macos"
\end{lstlisting}

Want to see what time it is?

\begin{lstlisting}[language=Lattice, numbers=none]
lattice> time()
=> 1740000000000
lattice> time_format(time(), "%Y-%m-%d %H:%M:%S")
=> "2025-02-19 16:00:00"
\end{lstlisting}

Need to know your working directory?

\begin{lstlisting}[language=Lattice, numbers=none]
lattice> cwd()
=> "/Users/alex/projects/lattice"
\end{lstlisting}

Every built-in function is available in the REPL, and tab completion helps you find them.

\subsection{Error Recovery}\label{subsec:error-recovery}
\index{REPL!error recovery}

When you make a mistake, the REPL prints an error and keeps going.
Your session is not lost:

\begin{lstlisting}[language=Lattice, numbers=none]
lattice> let x = 42
lattice> x.push(3)
error: push() requires an array, got Int
lattice> x
=> 42
\end{lstlisting}

The error did not corrupt the REPL state.
\lstinline|x| is still \lstinline|42|.
You can fix your mistake and try again.

This resilience comes from the REPL's error handling code in \filename{src/main.c}: when execution fails, the VM resets its stack and handler state, but globals and definitions are preserved.

\subsection{The Debugging REPL}\label{subsec:debug-repl}
\index{REPL!debugging}\index{breakpoint()}

For an even more powerful interactive experience, you can embed \lstinline|breakpoint()|\index{breakpoint()} calls in your source files.
When the interpreter hits a breakpoint, it drops into an interactive REPL where you can inspect local variables, evaluate expressions in the current scope, and examine the call stack.

\begin{lstlisting}[language=Lattice, caption={Using breakpoint() for in-context debugging}]
fn process_order(items: [String], total: Float) {
    let discount = if total > 100.0 { 0.1 } else { 0.0 }
    let final_price = total * (1.0 - discount)
    breakpoint()  // drops into REPL here
    print("Final price: ${to_string(final_price)}")
}

process_order(["widget", "gadget"], 150.0)
\end{lstlisting}

When the breakpoint is hit, you'll have access to \lstinline|items|, \lstinline|total|, \lstinline|discount|, and \lstinline|final_price| right there in the REPL\@.
We'll cover the full debugger in Chapter~32 (\emph{Debugging}).

\subsection{Alternative REPL Backends}\label{subsec:alt-repls}
\index{REPL!backends}

By default, the REPL runs on the stack-based bytecode VM\@.
You can switch to the other backends:

\begin{lstlisting}[language=bash, numbers=none]
./clat --tree-walk    # tree-walking interpreter
./clat --regvm        # register-based VM
\end{lstlisting}

The tree-walking REPL can be useful if you encounter a bytecode compiler bug and need a fallback.
The register VM REPL is more experimental.
In normal use, the default bytecode VM REPL is the right choice.

\begin{warnbox}{The Self-Hosted REPL}
Lattice also includes a self-hosted REPL written in Lattice itself (\filename{repl.lat}), built on the \lstinline|is_complete|, \lstinline|lat_eval|, and \lstinline|input| built-ins.
You can run it with \lstinline|./clat repl.lat|.
It lacks tab completion and readline integration, but it's an interesting demonstration of the language's metaprogramming capabilities.
The C-based REPL described in this chapter is what you should use day-to-day.
\end{warnbox}

\subsection{A Complete Exploratory Session}\label{subsec:complete-session}

Let's tie everything together with a realistic REPL session.
Imagine you are exploring Lattice's array methods for the first time:

\begin{lstlisting}[language=Lattice, numbers=none]
lattice> let scores = [85, 92, 78, 95, 88, 72, 91]
lattice> scores.len()
=> 7
lattice> scores.sort()
=> [72, 78, 85, 88, 91, 92, 95]
lattice> scores.min()
=> 72
lattice> scores.max()
=> 95
lattice> scores.sum()
=> 601
lattice> scores.filter(|s| s >= 90)
=> [92, 95, 91]
lattice> scores.filter(|s| s >= 90).len()
=> 3
lattice> let avg = to_float(scores.sum()) / to_float(scores.len())
lattice> avg
=> 85.85714285714286
lattice> scores.map(|s| if s >= 90 { "A" } else { "B" })
=> ["B", "A", "B", "A", "B", "B", "A"]
lattice> scores.enumerate()
=> [[0, 85], [1, 92], [2, 78], [3, 95], [4, 88], [5, 72], [6, 91]]
\end{lstlisting}

In under a minute, you've explored sorting, filtering, mapping, aggregation, and enumeration---all without writing a file.
That is the power of an interactive session.

%% ───────────────────────────────────────────────────────────
\section{Exercises}\label{sec:ch02-exercises}

\begin{enumerate}
  \item \textbf{REPL Warm-Up.}
  Launch the REPL and experiment with arithmetic: try \lstinline|2 ** 10| (does it work? what does Lattice use for exponentiation?), integer vs.\ float division (\lstinline|7 / 2| vs.\ \lstinline|7.0 / 2.0|), and the modulo operator.

  \item \textbf{Define and Iterate.}
  In the REPL, define a function \lstinline|fn is_palindrome(s: String) -> Bool| that checks whether a string reads the same forwards and backwards.
  Hint: \lstinline|s.reverse()|.
  Test it with \lstinline|"racecar"|, \lstinline|"hello"|, and \lstinline|"madam"|.

  \item \textbf{Struct Exploration.}
  Define a \lstinline|struct Color| with fields \lstinline|r: Int|, \lstinline|g: Int|, \lstinline|b: Int| in the REPL\@.
  Create several instances.
  Then define a function \lstinline|fn to_hex(c: Any) -> String| that converts a Color to a hex string like \lstinline|"#FF8800"|.
  Hint: use \lstinline|format()| or string concatenation.

  \item \textbf{Tab Completion Safari.}
  In the REPL, type \lstinline|str| and press \texttt{Tab}.
  What completions appear?
  Type \lstinline|"hello".| and press \texttt{Tab} twice.
  How many string methods are available?
  Try \lstinline|Map::| and \texttt{Tab}.

  \item \textbf{The Mystery Function.}
  Without looking it up, use the REPL to figure out what \lstinline|ord("A")| returns.
  Then figure out what \lstinline|chr(72)| returns.
  Can you write a one-liner that converts a string to an array of character codes?
\end{enumerate}

%% ───────────────────────────────────────────────────────────
\section*{What's Next}\label{sec:ch02-whats-next}

Now that you know how to use the REPL as your personal laboratory, it's time to learn what you can put into it.
In the next chapter, we will explore Lattice's value types---integers, floats, booleans, strings, nil, and unit---and the operators that work on them.
We'll see how Lattice thinks about data at the most fundamental level: the shape of things.
